\documentclass[10pt,a4paper]{amsart}
\usepackage[margin=19mm]{geometry}
\usepackage{amsmath,amssymb,mathtools}
\usepackage[scr=rsfs,cal=euler]{mathalfa}
\usepackage{xfrac}
\usepackage[unicode,hidelinks,bookmarks=false]{hyperref}
\newcommand{\ie}{i.e.\ }
\newcommand{\cf}{cf.\ }
\newcommand{\resp}{respectively\ }
\newcommand{\Subsection}{\subsection}
\providecommand{\nobreakdash}{\nobreak\hbox{-}}
\setlength{\emergencystretch}{2em}
\multlinegap0pt
\usepackage{tikz}
\usetikzlibrary{cd}
\usepackage{amssymb}
\usepackage{mathtools}
\usepackage{xfrac}
\usepackage{float}
\usepackage[shortlabels]{enumitem}
\usepackage{tikz}
\newtheorem{lemma}{Lemma}
\newtheorem{proposition}{Proposition}
\usepackage{cleveref}
\crefname{lemma}{Lemma}{Lemmas}
\crefname{proposition}{Proposition}{Propositions}
\newcommand{\Q}{\mathbb{Q}}
\newcommand{\Z}{\mathbb{Z}}
\newcommand{\lmfdb}[1]{%
  \href{https://www.lmfdb.org/EllipticCurve/Q/#1}{\texttt{#1}}%
}
\newcommand\sm[1]{\begin{bsmallmatrix}#1\end{bsmallmatrix}}
\title{Torsion of $\Q$-curves over number fields of small odd prime degree}
\author{Ivan Novak}
\date{Source: arXiv:2506.00753v2 (CC BY 4.0)}
\begin{document}
\maketitle

\noindent\emph{Source and licence.} Torsion of $\Q$-curves over number fields of small odd prime degree. Version arXiv:2506.00753v2 (CC BY 4.0), distributed by arXiv under the Creative Commons Attribution 4.0 International licence, http://creativecommons.org/licenses/by/4.0/, as recorded in the arXiv licence metadata for this exact version and shown on the abstract page https://arxiv.org/abs/2506.00753v2. Full licence notice: \texttt{licenses/arxiv-2506.00753-CC-BY-4.0.txt}.\par\smallskip

\noindent\textit{Editorial scope.} Counted unit: the article's proposition eliminiraj15 (source0.tex lines 468--470). The article's complete original proof of it is delivered with it (source0.tex lines 471--509, one \texttt{proof} environment of 39 lines). No mathematics is written, repaired or re-derived: the statement and the proof are the article's own lines, in the article's own notation, and no retained line is substituted or re-flowed.  Retained uncounted as the source-local context the proof needs: lines 314--316; lines 442--444.  Nothing was omitted from any retained span, so the package carries no omission record.  No retained line contains a citation, so the wrapper carries no bibliography and no bibliography entry.  The retained lines contain the article's own diagram code, which the wrapper typesets with the standard tikz packages; no image file is delivered and the package declares no asset dependency.  Editorial formatting only, in addition: an AMS \texttt{amsart} preamble that restates the article's own declarations and macros used by the retained lines, namely the environments \texttt{lemma}, \texttt{proposition} on separate counters, together with the standard packages those lines need.  The retained environments print in this document as Lemma 1, Proposition 1 in order of appearance, whereas the article prints the same items with the numbering of its own full document.  The complete native source of the article as distributed in the arXiv e-print for this exact version is bundled unchanged as \texttt{sources/179110/source0.tex}.
\begin{lemma}\label{twistanje}
    Let $E$ be an elliptic curve over a number field $K$. Let $E'$ be a quadratic twist of $E$ by $d \in K^\times$. Let $\psi:G_K \to \{-1, 1\}$ be the quadratic character defined by $\sigma(\sqrt{d})=\psi(\sqrt d) \sqrt d$ for $\sigma \in G_K$. Denote by $\rho_{E, N}$ and $\rho_{E', N}$ the mod $N$ Galois representations of $E$ and $E'$ respectively. Then, for an appropriate choice of bases of $E[N]$ and $E'[N]$, we have $$\rho_{E, N}(\sigma)=\psi(\sigma)\rho_{E'N}(\sigma).$$
\end{lemma}

\begin{proposition}\label{eliminiraj11}
    The group $\Z/11\Z$ does not occur as a torsion subgroup of a $\Q$-curve over a cubic number field.
\end{proposition}

\subsection{Eliminating $\Z/15\Z$} \begin{proposition}\label{eliminiraj15}
    The group $\Z/15\Z$ does not occur as a torsion subgroup of a $\Q$-curve over a cubic number field.
\end{proposition}

\begin{proof}
Suppose that some $\Q$-curve $E$ over a cubic field $K$ has a point $P$ of order $15$.
We have the following situation: 
\begin{center}
\begin{tikzcd}
    E/K \arrow[r, rightarrow, "\phi"] & E'/K \arrow[r, rightarrow, "\sim"] & E''/\Q.
\end{tikzcd}
\end{center}
Here $E'$ is an elliptic curve over $K$ with rational $j$-invariant, and $E''$ is a quadratic twist of $E'$ which is defined over $\Q$. 

Then, similarly as in the proof of Proposition \ref{eliminiraj11}, we can conclude that $E''$ has a $3$-isogeny and a $5$-isogeny over $\Q$, so it must have a $15$-isogeny over $\Q$. There are four possible $j$-invariants for which this happens.

As in the proof of Proposition \ref{eliminiraj11}, we can take one specific elliptic curve for each of the $j$-invariants.

We will look at curves from the isogeny class \lmfdb{1600.i} on LMFDB. Note that $-I$ is contained in the mod $15$ image.

By \Cref{twistanje}, we have $\rho_{E', 15} \sim \psi \cdot \rho_{E'', 15}$, where $\psi$ is some quadratic character. Furthermore, any element of $\rho_{E', 15}(G_K)$ has $1$ as an eigenvalue when reduced modulo $3$ and modulo $5$. This means that any element of $\rho_{E'', 15}(G_K)$ either has $1$ as an eigenvalue modulo $3$ and modulo $5$, or it has $-1$ as an eigenvalue modulo $3$ and modulo $5$. Now, for each of the four curves $E''$ from the class 1600.i, we do the following:

\begin{itemize}
    \item List all index $3$ subgroups $H$ of $\rho_{E'', 15}(G_\Q)$.
    \item In each of the subgroups $H$ from the list, find a matrix $A$ such that at least one matrix among $A \pmod 3$,  $A \pmod 5$ doesn't have $1$ as an eigenvalue, and at least one matrix among $A \pmod 3$, $A \pmod 5$ doesn't have $-1$ as an eigenvalue. 
\end{itemize}

The existence a matrix $A \in H$ with this property implies that $H$ cannot be the image $\rho_{E'', 15}(G_K)$.

\begin{enumerate}
    \item \textbf{Case 1600.i1} There are three index $3$ subgroups, so we have three cases. They are listed in order in which they appear in the auxiliary Magma program.
    \begin{enumerate}
        \item The matrices $\sm{2 & 5 \\ 0 & 1}$ and $\sm{1 & 10 \\ 0 & 11}$ are in the image $\rho_{E'', 15}(G_K)$. They both have to be in $\rho_{E', 15}(G_K)$ as their negatives don't have $1$ as an eigenvalue mod $5$. However, their product doesn't have $1$ as an eigenvalue mod $3$. 
        \item The matrices $\sm{12 & 1 \\ 10 & 6}$ and $\sm{6 & 5 \\ 5 & 6}$ are in $\rho_{E'', 15}(G_K)$. They don't have $-1$ as eigenvalues mod $5$, so they themselves have to be in $\rho_{E', 15}(G_K)$. However, their product equals $-I$ mod $3$, which doesn't have $1$ as an eigenvalue mod $3$, so we reach a contradiction. 
        \item Similarly as in the previous case, we end up with $-I$ in $\rho_{E', 15}(G_K)$, which is a contradiction.
    \end{enumerate}
    
    \item \textbf{Case 1600.i2}
    Each of the index $3$ subgroups contains at least one among $\sm{2 & 0 \\ 0 & 11}$ and $\sm{2 & 9 \\ 0 & 11}$. 
    \item \textbf{Case 1600.i3} Each of the index $3$ subgroups contains a matrix of the form $\sm{11 & * \\ 0 & 2}$.
    \item \textbf{Case 1600.i4} The matrix $\sm{11 & 0 \\ 0 & 11}$ is in all index $3$ subgroups.
\end{enumerate}
\end{proof}

\end{document}
